%% =================================================================
%% Wei, Mei, Zhao, Xiao, Sun, Zhang, Guo.
%% "Nondestructively identifying the mechanical behavior of soft
%%  tissues using surface deformation with an explicit inverse
%%  approach."
%% Appl. Math. Modelling 134 (2024) 126--147.
%%
%% Reconstruction of page 10 (printed p. 135) as study notes.
%% This page is mirror-structured to page 9: two big figures
%% (Fig. 9 for Sample 2, Fig. 10 for Sample 3) comparing Explicit
%% vs Implicit methods with 3, 6, and 9 displacement datasets,
%% followed by minimal prose.
%% Prose: extracted verbatim from the PDF text layer via pdftotext.
%% Math:  none on this page.
%% Figures 9 & 10: placeholder boxes with captions + visual descriptions.
%% =================================================================

\documentclass[11pt]{article}

\usepackage[margin=1in]{geometry}
\usepackage{amsmath, amssymb}
\usepackage{mathrsfs}
\usepackage{bm}
\usepackage{xcolor}
\usepackage{fancyhdr}

\pagestyle{fancy}
\fancyhf{}
\lhead{\small Y.~Mei et al.}
\rhead{\small Applied Mathematical Modelling 134 (2024) 126--147}
\cfoot{\small 135 $\to$ \arabic{page}}
\renewcommand{\headrulewidth}{0.4pt}

\setcounter{page}{1}

\begin{document}

% ---------- Figure 9 placeholder ----------
\noindent
\begin{center}
\fbox{\parbox{0.92\textwidth}{\centering\vspace{1em}
\textbf{[Figure 9 — placeholder]}\\[0.8em]
Six-panel side-by-side comparison of the \textbf{Explicit Inverse
Method} (left column) vs.\ \textbf{Implicit Inverse Method} (right
column) on \textbf{Sample 2} (the two-layer target from Fig.~6b,
nominal SM values $\mu^0 \approx 3$\,MPa, $\mu^1 \approx 10$\,MPa),
using 3, 6, and 9 surface displacement datasets as inputs.\\[0.4em]
\textbf{Explicit (left column)}: bilayer structure recovered with a
clean interface. Colorbar ranges:
(a) 3 fields, SM $\in [1.00, 11.49]$\,MPa;
(b) 6 fields, SM $\in [1.00, 11.63]$\,MPa;
(c) 9 fields, SM $\in [1.00, 11.47]$\,MPa.\\[0.3em]
\textbf{Implicit (right column)}: bilayer structure recovered but
with a visibly rougher interface between the two layers; more
artifacts at the top edge. Colorbar ranges:
(d) 3 fields, SM $\in [1.00, 10.95]$\,MPa;
(e) 6 fields, SM $\in [1.00, 10.77]$\,MPa;
(f) 9 fields, SM $\in [1.00, 11.13]$\,MPa.\\[0.3em]
\textbf{Takeaway}: both methods recover the two-layer structure;
the explicit method produces a cleaner interface, while the implicit
method's interface has more small-scale roughness. Slight overshoot
in the top-layer SM values (colorbar max above the nominal 10\,MPa)
is visible in all six panels.
\vspace{1em}}}
\end{center}
\vspace{0.3em}
\begin{center}
\textbf{Fig. 9.} The reconstructed results with varying 3, 6 and 9
surface displacement datasets corresponding to Sample 2 (Without
noise, Unit: MPa).
\end{center}

\vspace{0.8em}

% ---------- Figure 10 placeholder ----------
\noindent
\begin{center}
\fbox{\parbox{0.92\textwidth}{\centering\vspace{1em}
\textbf{[Figure 10 — placeholder]}\\[0.8em]
Six-panel side-by-side comparison of the \textbf{Explicit Inverse
Method} (left column) vs.\ \textbf{Implicit Inverse Method} (right
column) on \textbf{Sample 3} (the three-layer target from Fig.~6c:
blue bottom layer $\mu^0 \approx 3$\,MPa, green middle layer
$\mu^1 \approx 6$\,MPa, red top layer $\mu^2 \approx 10$\,MPa),
using 3, 6, and 9 surface displacement datasets as inputs.\\[0.4em]
\textbf{Explicit (left column)}: three-layer stacked structure
recovered, with the middle (green) layer visibly distinct from the
top (red) and bottom (blue) layers. Colorbar ranges:
(a) 3 fields, SM $\in [1.00, 10.49]$\,MPa;
(b) 6 fields, SM $\in [1.00, 10.35]$\,MPa;
(c) 9 fields, SM $\in [1.00, 10.44]$\,MPa.\\[0.3em]
\textbf{Implicit (right column)}: three-layer structure recovered
but with fuzzier interfaces between adjacent layers. Colorbar
ranges:
(d) 3 fields, SM $\in [1.00, 10.11]$\,MPa;
(e) 6 fields, SM $\in [1.00, 10.26]$\,MPa;
(f) 9 fields, SM $\in [1.00, 10.19]$\,MPa.\\[0.3em]
\textbf{Takeaway}: as layer count increases (1 $\to$ 2 $\to$ 3),
the quality gap between explicit and implicit widens. Sample 3's
three-layer structure is harder for the implicit method to resolve
cleanly than Sample 2's two-layer structure. More displacement
datasets (3 $\to$ 6 $\to$ 9) improve both methods but do not close
the gap.
\vspace{1em}}}
\end{center}
\vspace{0.3em}
\begin{center}
\textbf{Fig. 10.} The reconstructed results with 3, 6 and 9 surface
displacement datasets corresponding to Sample 3 (Without noise,
Unit: MPa).
\end{center}

\vspace{1em}

% ---------- Continuation prose (from page 9) + 5% noise paragraph ----------
\noindent
slightly larger than that for the explicit inverse method (see
Table~1). Furthermore, it's worth noting that the quality of the
reconstruction achieved by both inverse methods improved with an
increase in the total number of measurements. By utilizing 9
boundary displacement measurements, the average relative error
could be reduced to less than 10\,\%. However, it should be
mentioned that each layer was not accurately mapped as the
homogeneous case when using the implicit inverse method.

Additionally, we increased the noise level to 5\,\% and observed
that the reconstructed results obtained by the explicit inverse
method still outperform those obtained by the implicit inverse
method (see Fig.~14, Fig.~15, Fig.~16). For the explicit inverse
method, the resulting relative error in the case with 5\,\% noise
is generally larger than that in the case with 1\,\% noise (see
Table~1). However, even with a 5\,\% noise level, the relative
error for the explicit inverse method remains lower than that for
the implicit inverse method

% [page ends here — continues on page 11]

\end{document}
